
%----------------------------------------------------------------------------------------
%	PACKAGES AND OTHER DOCUMENT CONFIGURATIONS
%----------------------------------------------------------------------------------------

\documentclass[12pt]{article} % Default font size is 12pt, it can be changed here
\usepackage{hyperref}
\usepackage{amsmath}
\usepackage{geometry} % Required to change the page size to A4
\geometry{a4paper} % Set the page size to be A4 as opposed to the default US Letter

\usepackage{graphicx} % Required for including pictures

\usepackage{float} % Allows putting an [H] in \begin{figure} to specify the exact location of the figure
\usepackage{wrapfig} % Allows in-line images such as the example fish picture

\usepackage{lipsum} % Used for inserting dummy 'Lorem ipsum' text into the template

\linespread{1.2} % Line spacing

%\setlength\parindent{0pt} % Uncomment to remove all indentation from paragraphs

\graphicspath{{./Pictures/}} % Specifies the directory where pictures are stored
\begin{document}

%----------------------------------------------------------------------------------------
%	TITLE PAGE
%----------------------------------------------------------------------------------------

\begin{titlepage}

\newcommand{\HRule}{\rule{\linewidth}{0.5mm}} % Defines a new command for the horizontal lines,
%change thickness here
\center % Center everything on the page
\includegraphics[width=\textwidth]{Glasgow}\\[1.5cm]
\textsc{\LARGE Algorithmics}\\[0.5cm] % Major heading such as course name
\textsc{\Large Assessed Exercise}\\[0.5cm] % Minor heading such as course title

\HRule \\[0.4cm]
{ \huge \bfseries Suffix Tree Applications}\\[0.4cm]
\\
\emph{State \& Strategy}
\HRule \\[1.5cm]

\begin{minipage}{0.4\textwidth}
\begin{flushleft} \large
\emph{Author:}\\
Garry \textsc{Sharp}\\
0801585s\\ % Your name
\end{flushleft}
\end{minipage}
~
\begin{minipage}{0.4\textwidth}
\begin{flushright} \large
\emph{Supervisors:} \\
Dr. D. \textsc{Manlove}\\ % Supervisor's Name
\end{flushright}
\end{minipage}\\[4cm]

{\large \today}\\[3cm] % Date, change the \today to a set date if you want to be precise

\vfill % Fill the rest of the page with whitespace

\end{titlepage}

\tableofcontents
\newpage
\section{Task 1 (Search Suffix Tree)}
The task here is implement a method searchSuffixTree in the class SuffixTreeAppl
which searches over a global instance of a SuffixTree known as t with the search
term being an array of bytes x.
\subsection{State}
My solution works in an effective manner to search the tree in O(n) time. It should be
noted that the solutions to Task \ref{task2}, Searching for all occurrences relies heavily on
this method. This means that there is a very small amount(almost negligable) amount of redundant
code when searching (sets a global variable). 
\subsection{Solution}
Firstly, lets examine the variables that are used. I keep track of an x index (initialised
to 0), a node index (index of the character being compared at the node), a boolean called match
(which is set to true/false depending on whether there is a match), an integer startLocation
(initialised to -1) as well as a SuffixTreeNode called currentNode (initialised to the child of the
root of the tree).

My solution compares the relevant values of x with the relevant values held at the node. If all
characters are matched successfully (and there are more in subsequent nodes) then I perform the
same check at the child of that node. If all characters have been matched and no more are left then
the loop terminates.

If there is a mismatch at any characters then I examine the sibling that node until, either there
is another match (is a mismatch occurs at that node it looks at the sibling, otherwise if fully
matched goes to the child) or a mismatch at which point the loop exits.

A variable t2Node is set throughout and this is so that Task 2 can be completed by recycling the
code here in Task 1. This node is eventually set to the final node that contains leafs of the
instances of the search term.

If a match does occur and startLocation is not already set, a variable startLocation is set to it's
left label of the currentNode. Otherwise is a mismatch occurs startLabel is reset (set to -1).

The program terminates by returning a Task1Info object (with position startLocation which will be
-1 if not found) to the main method which then displays it accordingly. The search executes in O(n)
time. This solution uses no external helper methods.
\newpage
\section{Task 2 (All Occurrences)}
The task here is very similar to Task 1, but with the result being equal to every instance of the
search term in the text (as opposed to just a single occurrence), specifically the starting indices
of every occurrence.
\label{task2}
\subsection{State}
This works in an effective manner by reusing the search functionality implemented in Task 1. This
is done by calling the searchSuffixTree method and then using the t2Node (which is set to the
branch node from which all decendent leaf nodes contain the occurrences). The solution executes in
O(n) time (plus build time).
\subsection{Solution}
As mentioned above, searchSuffixTree is called and t2Node is assigned accordingly. From here on I
use a recursive method on the child of the node until there are no more siblings. I basically look
at all nodes, where if there is no child (ie. node is a leaf) I add the value to a global variable
called occurrences, otherwise (if leaf) I recurse down the child of the currentnode and after move
to the sibling.

Finally I iterate through the value of occurences and add each of the suffixes of each node to the
Task2Info object which I then return. The recursive call to iterate through the values is
getLeafDecendants.
\newpage
\section{Task 3 (Longest Repeating Substring)}
\label{sec:t3}
The task here is to look through a tree and find the longest repeating substring, or to put it
another way, the longest multiple occurring sequence of characters. A repeating substring exists if
a node has two leaf nodes that are it's children (or siblings of the child in this implementation).
To get the LONGEST repeating substring, the node with the the greatest path depth must be
found. This is the value of the highest ancestor's (first item in the list's) leftLabel subtracted
from the last node in the list's (currentNode's) right label (provided that currentNode is valid ie.
has 2 leaf children). 
\subsection{State}
The solution detailed works as expected when tested over the data sets provided. 
\subsection{Solution}
The solution consists of several parts which I will detail individually. Firstly I should state
that my solution is not ideal as it does not visit each node only once, however, in compromise, the
readability offered by this is greatly improved. 

As explained above there are several schools of thought in calculating this, firstly there must be
a way of checking if a node is valid (ie. has at least 2 children that are leaf nodes). To aid
readability this is done in the method isValidBranch which takes a SuffixTreeNode as an argument
and returns true or false if it is "valid" (has at least 2 leaf node children).

In addition there must be some form of iteration over the tree in order to find the greatest path
depth, this is aided by a method called next. Next takes a list of SuffixTreeNodes as an argument
which represent a path from the highest ancestor (a child of the root of tree t) to the currentNode
(the furthest down descendant). Next determines what the next path should be and will either visit a
child (adding to the list and returning it), visit siblings (removing the last element of the list
and adding the sibling in it's place). Finally if all children/siblings have been visited then the last element is
removed continually until either the last node in the list has a sibling OR there is no more nodes in the list (meaning all paths have been visited).  The requirement of a list is because nodes do not have a .getParent() method. A local variable "visited" is featured as to
prevent an infinite loop occurring by removing a node from the list and then immidiately adding it
again, this is also used to allow the path to shrink by more than one node if necessary (ie. the immediate parent has no sibling). This is perhaps the most complex part of the solution (building the next method) and is what
consumed the majority of my time on this specific task.

Finally a method that takes a path (explained below) is needed, which takes the current path and
calculates it's depth by summing the individual lengths (rightlabel - leftlabel + 1).

The final solution basically combines all of the above methods, iterating through the tree until
the path equals null (terminating the iteration as all paths have been examined), and comparing it
the length of the deepest path (called bestLsrNodes in solution). if the currentPath is deeper than
the current best and also has a "valid node" (node with at least 2 child leafs) as the last entry,
then bestLsrNodes is updated accordingly.

Finally a Task3Info object is returned which contains the length of bestLsrNodes and two of the
suffixes (or positions) of the last node in the bestLsrNodes list.

NB. Because I calulate the length of the path every time a new path is found you may consider the
code inefficient, however, this was done intentionally for a number of reasons. Firstly, it aids in
the logical reading of the code, improving readability and description. Secondly given the time
contraints, I did not want to dwell for to long on a minor improvement in efficiency (takes
approximately 350ms to execute as it is currently). The improvement could be realised by the creation
of a bean/inner class which consists of a SuffixTreeNode and a depth (which is calculated in the next method).
This inner class could then be used in the path/bestLrsNodes instead of just the nodes, removing the need for
computing the depth later as it is stored along with the node. This was not done as the project description advices against using inner classes in the implementation of the SuffixTreeAppl class. The program still executes quickly (faster than construction) despite this.
\newpage
\section{Task4 (Longest Common Substring)}
The task here is to output the longest common substring over two files (ie. the longest sequence of
characters that is present in both some file a and also some file b.
\subsection{State}
The code (I think) works well as it does output what I think is the longest common string. That
said, when I have compared my results to the results of peers they have been different upon
occation, the longest common substring that is returned by comparing text1.txt and text2.txt is the same as the output on the handout sheet under "example output".
\subsection{Solution}
The code used here relies heavily upon the code detailed in Section \ref{sec:t3}.
It is therefore advisable that you are familiar with this method before reading this as I will
presume familiarity. 

Firstly, in order to use this method, a constructor for two files must be implemented in the
SuffixTree class. This is done by reading in the bytes to the first file, appending a \# symbol,
appending the bytes for the second file, and finally appending a \$ symbol. Variables such as
stringLen must here be set to the length of the total bytes -1 (as to not include the \$ symbol as a
character in the actual tree). No checking that either files contain a \$ or \# symbol is done as this is stated as a premumption in the constructor's javadoc. The tree is then built.

For the description below, file a and file b reference the first and second files' bytes that make
up part of the tree respectively.

The code is almost identical to the solution as detailed in task 3, however, there is a difference
in the sense that an additional check must be made AFTER checking is the node is valid (ie.
contains at least two leaf nodes) and BEFORE setting the bestLsrNodes to the path if the depth is
greater than what has previously been encountered. This step is simply making sure that AT LEAST
ONE of the suffixes of the last node in the path is less than the length of the first file (ie. it
is present in file a) and also that AT LEAST ONE OTHER suffix is greater than the length of the
first file (ie. it is present in file b). If both of these conditions are met, then and only
then is the bestLrsNodes updated, if this condition is not met then the two substrings are not
common, ie. It is not present in BOTH files. 

This solution suffers from the same minor flaw as task 3 in that it visits nodes multiple times to
calculate the path depth. Other than this it appears to work as expected.
\newpage

%----------------------------------------------------------------------------------------
%	BIBLIOGRAPHY
%----------------------------------------------------------------------------------------

%\begin{thebibliography}{99} % Bibliography - this is intentionally simple in this template

%\bibitem[Figueredo and Wolf, 2009]{Figueredo:2009dg}
%Figueredo, A.~J. and Wolf, P. S.~A. (2009).
%\newblock Assortative pairing and life history strategy - a cross-cultural
%  study.
%\newblock {\em Human Nature}, 20:317--330.
 
%\end{thebibliography}

%----------------------------------------------------------------------------------------

\end{document}
